\chapter{Introduction}
\label{chapter:Introduction}

\pagenumbering{arabic} % Sets page numbering style to arabic
\setcounter{page}{1} % Resets page counter to 1
\pagestyle{plain}


%\guidelines{{\bf Recommended Total Length of Introduction: About 5 pages.}}


%\section*{Introductory Paragraphs}
%\guidelines{{\bf Recommended Length: half a page.}}
%\guidelines{Please comment the heading of this section. Just start the chapter with the introductory paragraphs.}


%\guidelines{Motivate your work by writing about two paragraphs. This should be more or less an extended version of your abstract where you justify your statements using reference to the literature. }

Technical indicators play a central role in rule-based and algorithmic trading, where they are used to extract information about trends, momentum, and market participation from historical price and volume data. Prior research suggests that the profitability of individual technical trading rules is mixed and often varies across markets and time periods \citep{lento2007profitability}. This has motivated
continued interest in more structured trading systems that combine several
indicators rather than relying on a single rule
\citep{keshavarz2022trading}. At the same time, machine-learning-based
approaches have increasingly been proposed as a way to improve signal
selection in trading strategies \citep{sukma2025enhancing}. However, more complex models may also increase the risk of overfitting, which makes it relevant to investigate whether simpler and more interpretable models can improve robustness and generalisation without introducing unnecessary complexity
\citep{dietterich1995overfitting}.

This thesis addresses that issue by studying whether the confidence produced by
a decision tree can be used to distinguish between lower-quality and
higher-quality signals generated by a rule-based trading strategy on the gold
market. Rather than replacing rule-based logic with a fully data-driven model,
the study examines whether a simple and transparent machine-learning layer can
be used to evaluate candidate signals in a way that preserves interpretability
while supporting stronger out-of-sample performance. To investigate this, a
multi-indicator strategy based on EMA, RSI, MACD, DMI, KST, and MFI is
developed using historical H1 data for the broker symbol \texttt{XAUUSD} over
the period January 2022 to December 2024. Candidate trade signals are then
filtered by a decision tree trained on indicator-derived features, after which
the final strategy is evaluated on a held-out out-of-sample segment using a
deterministic Python backtesting pipeline built on MetaTrader~5 market data.

\section{Background}
%\guidelines{{\bf Recommended Length: between half a page and one page.}}
%\guidelines{Feel free to change the title of the section to something more specific, such as `Neural Networks'}
%\guidelines{Discuss background information related to the techniques to be explored in the work. }

Rule-based trading refers to trading systems in which trading decisions are generated from predefined rules based on historical market data \citep{ozturk2016heuristic}.

In many trading systems, several indicators are combined rather than used in isolation, since different indicators may capture different dimensions of market behaviour \citep{sukma2025enhancing}. In this study, this forms the basis of the initial trading strategy, which is built as a multi-indicator rule-based system applied to the gold market through the broker symbol XAUUSD.

Machine learning can be applied to such systems not only for prediction but also for signal filtering. In that setting, signals generated by a rule-based strategy are evaluated by a classification model and either accepted or rejected based on their characteristics \citep{wu2006decisiontree}. In this thesis, that role is assigned to a decision tree. Decision trees are relevant in this context because they express decisions as explicit if-then rules, which makes them comparatively easy to interpret \citep{nugroho2014decision}.

Because financial market data are noisy and time-dependent, model complexity must be handled carefully. A model that is too flexible may adapt to historical noise rather than patterns that generalize to unseen data, which makes overfitting an important concern in trading research \citep{dietterich1995overfitting}. 

\section{Problem}
%\guidelines{\bf Recommended Length: one page}
%\guidelines{Here you should identify a `general gap' in the literature. This should be something of general interest. Example: `Although problem X can be used as a basis to solve a large number of practical computational problems, problem X itself has been shown to be computationally hard. It is therefore relevant to explore new techniques to address problem X.}

Several studies have investigated rule-based trading strategies and technical indicators, yet empirical findings regarding their profitability remain mixed (e.g., \cite{lento2007profitability, curcio1997technical}. To enhance such strategies, increasingly complex approaches have been proposed, including evolutionary algorithms, probabilistic reasoning models, and advanced machine learning techniques (e.g., \cite{potvin2004generating, chou2014rule, sukma2025enhancing}).

Although these methods may improve in-sample performance, greater model complexity increases the risk of overfitting, potentially compromising robustness and out-of-sample generalisation (\cite{dietterich1995overfitting, hawkins2004problem}). 

\begin{problemformulation}
\label{problemformulationOne}
This raises the question of whether simpler machine learning methods can enhance traditional rule-based trading strategies without introducing excessive complexity or sacrificing generalisation performance.

\end{problemformulation}



\section{Research Question}
%\guidelines{\bf Recommended Length: half a page}

%\guidelines{
%Here you should identify a general interest research question that can be seen as an approach towards dealing with the problem you have described above. 
%}

\begin{researchquestion}
\label{researchQuestion}

How does the confidence of a decision-tree signal filter relate to the out-of-sample performance of a rule-based trading strategy on the gold market?

\end{researchquestion} 

To address this question, the study first develops a rule-based trading strategy
based on technical indicators. Then a decision tree is applied to evaluate the
quality of the candidate trade signals, following the general approach of
\citet{nugroho2014decision}. For each signal, the model produces a confidence
score, interpreted as the predicted probability of a favourable realised outcome.
These confidence scores are divided into ordered ranges, and then the
performance of the strategy is evaluated in these confidence ranges. Since
financial prediction models are vulnerable to overfitting, analysis places
particular emphasis on out-of-sample evaluation and robust backtesting
procedures, as follows \citet{bailey2017pbo}. The study therefore isolates a
final held-out market segment for evaluation and compares the confidence ranges
under identical transaction-cost and exit assumptions.

%\guidelines{Make some comments about how you are going to address you research question. This is more or less a brief synopsis of the methodology section. }


\section{Delimitations}

%\guidelines{Specify delimitations related to your research question/research goal. Examples what markets are you considering,  what Market sessions, what timeframes, etc. } 

This study is limited to gold and therefore focuses exclusively on one commodity market. The findings cannot be directly generalised to other markets, such as equity indices, individual stocks, foreign exchange, or cryptocurrencies. In addition, the analysis is based on hourly data only, meaning that alternative trading horizons and higher-frequency strategies are excluded from the scope of the study.

The study is further delimited to a specific hybrid approach in which a decision tree model is used to evaluate candidate signals generated by a rule-based trading strategy. More advanced machine-learning models, such as neural networks or ensemble methods, are not examined. The purpose is to evaluate whether a simple and interpretable model can distinguish between weaker and stronger trade signals on the basis of its confidence estimates.
